\section{Hardware and scope: full account}
\label{sm:scope}

Three claims one could make for this method are larger than the quantity that fixes them: the hardware claim
is fixed by a subspace size, the chemistry claim by a set of sector dimensions, and the advantage claim
by the same bond dimension this paper puts forward as its own cost tracker. This section states each
with its fixing quantity attached, and derives---rather than concedes---the region in which this method
has no advantage over a classical one.

\paragraph{The two hardware runs.} Both device runs are preserved in full and they are the entirety of the quantum hardware used in this
work. The first was executed on the IBM~Heron processor \texttt{ibm\_fez} (job
\texttt{d9s16avpemts73ct6g8g}): the $L=6$ Hubbard chain at $U/t=4$, $\eta=0.15\,t$, $12$ qubits,
$3.5\times10^{5}$ computational-basis shots---$5\times10^{4}$ on each of seven circuits
$k=0,\dots,6$, of which $k=0$ is the $t=0$ reference---pooled into one subspace,
under TREX readout twirling, Pauli twirling and dynamical decoupling and
with no zero-noise extrapolation
(Fig.~\ref{fig:heron}(a) is the result; the circuits prepared the computational-basis $(N{+}1)$ determinant with orbitals $0$--$3\!\uparrow$ and $0$--$2\!\downarrow$ occupied rather than $\hat c^{\dagger}_p\lvert\Psi_0\rangle$, and their hopping layer applies $R_{xx}R_{yy}$ rotations without Jordan--Wigner strings, so they are not the primitive of Fig.~\ref{fig:circ}; because the recovered subspace is the whole sector, neither affects the reported $A_{\rm hw}$). Self-consistent configuration recovery returned $|\mathcal S|=300$ determinants of an
$(N{+}1)$ sector of dimension $300$; the recorded relative $L_1$ is $0.0$, and the deposited provenance
string says why in its own words---the reconstructed $A_{\rm hw}$ equals the exact $A$ byte for byte
because the recovered subspace spans the whole sector.

That number is exact \emph{by coverage}, and the distinction is not rhetorical: at
$|\mathcal S|=\dim$, Rayleigh--Ritz on $\operatorname{span}(\mathcal S)$ \emph{is} diagonalization of
the full sector, $P_{\mathcal S}HP_{\mathcal S}=H$, and a vanishing error is an identity rather than a
measurement of device fidelity. The run establishes that the pipeline executes end to end on a
processor---transpiled circuits, device counts, recovery to the correct particle-number sector, and the
classical Lehmann assembly of Sec.~\ref{sec:method}---which is what a proof of principle is, and it
establishes nothing about accuracy at incomplete coverage, the regime the rest of this paper is about.
The runs are therefore described as \emph{executed on} the named processors, not as demonstrating the method on them, and the coverage qualification applies wherever they are mentioned.

\paragraph{The second run, and a comparison the record cannot decide.} Stretched $\mathrm{N_2}$ ($R=2.0$~\AA,
CAS(10e,12o)/cc-pVDZ, $24$ qubits) was executed on \texttt{ibm\_marrakesh} (job
\texttt{da125f2ein7c73bcsqs0}, $10^{4}$ shots, double-factorized Trotter circuit). Averaged over $8$
recovery seeds on the cached device counts it reaches $0.59\pm0.14$~mHa of the exact active-space value
(best seed $0.40$~mHa), while the noiseless statevector simulation of the same shallow circuit stops at
$29.5$~mHa (Fig.~\ref{fig:heron}(b)). That contrast invites the reading that device noise \emph{helps}, and the record cannot support it, because the mechanism named for it is a statement about a quantity the
record does not contain: recovery is said to turn corrupted bitstrings into valid configurations and
thereby \emph{enlarge} the variational subspace---a valid-shot economy, not a fidelity gain---so the
comparison is decided by the two recovered subspace sizes, and neither is stored. The same gap runs
through the controlled sweep, whose script saves the size of the naively post-selected subspace and not
of the recovered one, so the two places where an error \emph{below} the noiseless value is
reported---CO from $1.73$ to $0.63$~mHa and $\mathrm{C_2}$ from $0.184$ to $0.026$~mHa, between
$\varepsilon=0$ and $\varepsilon=20$--$30\%$---rest on the same unrecorded variable, at the price of one
line in a script that already builds the subspace.

The mechanism itself is not in dispute. On the dissociation of
$\mathrm{N_2}$ in a large active space, with the same LUCJ ansatz used for the run above,
Vaquero-Sabater, Carreras, Broers, Shirakawa, Yunoki and Casanova report that ``sampling noise can
enhance Hilbert-space exploration by generating additional configurations beyond those supported by
the ideal ansatz,'' and that ``when combined with configuration recovery, this leads to systematically
improved energies''~\cite{vaquero2026}: the same molecule, the same ansatz and the same mechanism, and there it is carried by a controlled comparison. What this work cannot claim is its own measurement of the effect, not the effect: our two arms differ in
the one quantity neither of them recorded, so they cannot be evidence either way, whatever the
literature independently establishes.

What \emph{is} controlled is the comparison in which both arms consume the same corrupted bitstrings at
the same rate with the same seeds and differ only in the recovery step, and there the gain from
recovery is measured and grows with the noise: at $\varepsilon=2\%$ eighteen of nineteen molecules
reach chemical accuracy with recovery against seventeen without it, at $\varepsilon=30\%$ recovery
roughly halves the error on $\mathrm{N_2}$ ($31.9\to13.2$~mHa, Fig.~\ref{fig:noiserec}), and on the
Hubbard chain it holds the spectral relative $L_1$ near $0.05$ out to $\varepsilon\approx16\%$
(Fig.~\ref{fig:noise})---seed-averaged, with $\pm1$~s.d.\ bands, at matched noise. That is the noise
result of this work.

\paragraph{Scale.} These runs are $12$ and $24$ qubits, and the same class of observable has been
executed elsewhere at larger scale---retarded spin Green's functions of the one-dimensional
Fermi--Hubbard model on IBM~Heron~r2 at up to $30$ qubits~\cite{vilchez2025}, an ARPES-type
$A(k,\omega)$ on a $27$-site chain at $54$ qubits~\cite{granet2026akw}, dynamical structure factors
on trapped ions against scattering data~\cite{quantinuum_dsf2026}---so our device record is smaller
than every entry in that class and is not offered as a scale result. It is not the smallest such record
either, and we name the one below it: a spectral function was reconstructed on a trapped-ion device
from a four-qubit Hubbard dimer~\cite{greenediniz2024}. What is new in our run
is the object assembled, the $(N{\pm}1)$-sector Lehmann construction from a bitstring-sampled subspace,
at a size where the coverage question is degenerate, and Sec.~\ref{sec:prices} says how degenerate in
shots: the $3.5\times10^{5}$ spent---seven circuits at $5\times10^{4}$ shots each, pooled into one
subspace---are $14.4$ times the $(2.43\pm0.34)\times10^{4}$ that Table~\ref{tab:shots} measures for one
orbital seed and one branch at $L=6$, so complete coverage of a $300$-dimensional sector is what the
budget predicts rather than a surprise.\footnote{Table~\ref{tab:shots} replicates the repository protocol, $K=18$ with $19$ snapshots,
at $U/t=8$ rather than the $U/t=4$ of the device run; this is a consistency statement between two
budgets, not an identity.} The first size at which a device run would test the reconstruction rather
than the pipeline is $L=8$: sector $3\,920$, at $(2.70\pm0.13)\times10^{5}$ shots per orbital seed and
branch---a factor $1.3$ \emph{below} what was already spent here, so what confined this run to $L=6$
was not the shot budget. The momentum-resolved observable at that size is a different matter: sixteen
channels, ${\approx}4\times10^{6}$ shots, an order of magnitude above this job.

\paragraph{The region of no advantage, in full.} The resource statement of Sec.~\ref{sec:nogo} has a consequence for the standing of this method, and we
derive it here rather than wait to have it put to us. What tracks the determinant support is $\chi$,
the minimal bond dimension across a bipartition in the working basis (Sec.~\ref{sec:nogo:chi}). But
$\chi$ is, by definition, the cost parameter of a matrix-product representation of the same state. A
cost claim built on a basis-dependent boundary quantity therefore cannot select for the quantum side of
the comparison: wherever our own diagnostic says the determinant support is small it says a
matrix-product state is small too, and both methods are charged the same number. The thesis of
Sec.~\ref{sec:nogo} implies its own region of no advantage, and implies it more sharply than an
external objection would.

The instance is in the literature and it is not marginal. Ouyang, Chi and Chan reproduce the $7\,260$
trajectories of the Fermi--Hubbard dynamics experiment of Hartnett \emph{et al.}---$120$ qubits, up to
$90$ Trotter steps~\cite{qctrlfh2026}---by tensor contraction at bond dimension $32$, faster and to
higher accuracy than the device run~\cite{ouyang2026}: the
same pattern in which tensor-network and Pauli-propagation methods have overtaken earlier hardware
demonstrations, the first ``quantum utility'' claim~\cite{kim2023} among
them~\cite{tindall2024,begusic2024,kemper2026}, and in which unitary cluster-Jastrow sampling is
itself now classically reproduced, including by an extended matchgate
simulator~\cite{belagali2026,extraferm}---against which the standing worst-case hardness result for
that circuit family, its embedding of instantaneous-quantum-polynomial
sampling~\cite{bremner2016,hafid2025}, is not a contradiction but a statement about a different
regime. Two architecture-independent results point
the same way: polynomial-time classical estimation of expectation values in noisy
circuits~\cite{schuster2025}, and classical estimation in noiseless chaotic circuits across
geometries~\cite{angrisani2025}. On the ground-state energy the strong classical
selectors~\cite{holmeshci2016,sharmashci2017,tubmanasci2016} already win on
cost~\cite{reinholdt2025,zhangotten2026,gaberle2026}, and the case for a generic exponential advantage
in ground-state chemistry was found wanting before that~\cite{leechan2023}, and we concede it without reservation.

One framing does not survive. It is tempting to argue that the energy is where classical methods have caught up and to move the target to dynamics; in one dimension that inverts,
because dynamical DMRG computed all three of the momentum-resolved channels reported here nineteen
years ago~\cite{benthien2007} (Sec.~\ref{sec:intro}, Sec.~\ref{sec:physics}).
Dynamics in one dimension is where classical methods arrived first, and no framing of this work
says otherwise. Within the region this work measures---one-dimensional Hubbard rings at $L\le14$, the
nineteen active spaces, and a largest minimal bond dimension of $42$ over the $74$ deposited exact
ground states, attained on a molecular active space, the largest on the lattice being $38$ in the
two-leg ladder---there is no computation performed
here that a classical method cannot perform, and we make no advantage claim at any size reached.
Deriving one's own no-advantage region is a stronger position than conceding it; the credit does not
change the fact, and the fact is the one just stated.

The same derivation says where the boomerang stops, and that is the measurement this subsection closes
on. The controlled contrast of Sec.~\ref{sec:nogo:chi} is a search for a target on which the two costs
come apart: at one-body magic matched to within $1\%$, the two-leg ladder carries $\chi=13$ against the
chain's $5$ at $L=6$ and $\chi=38$ against $8$ at $L=8$, while its determinant support exceeds the
chain's by $1.52$, $1.68$ and $2.24\times$. The quantity to look for is a target whose minimal bond
dimension is large in \emph{every} ordering while the sampled support stays compact---a property of
geometry and orbital connectivity rather than of dimensionality, so two-dimensional lattices and
molecular orbital graphs are where it would be measured, and nothing here measures it.
Section~\ref{sec:frontier} supplies the second coordinate of that search: in the determinant basis the
Krylov space of the probe already has coordinate support on every symmetry-allowed determinant at $L=8$ ($3902$ of $3920$), so zero leakage at full captured weight requires a different working basis; whether a non-vacuous certificate at small sector fraction does is open (Sec.~\ref{sec:frontier:obstruction}).
